% Gomory-Hu Tree

A tree structure representing the minimum cuts (and maximum flows) between all pairs of vertices in an undirected graph using $N-1$ max-flow calls.

\begin{lstlisting}
struct GomoryHu {
    struct Edge { int u, v; long long w; };
    int n;
    vector<Edge> tree;

    GomoryHu(int _n) : n(_n) {}

    // Requires a Dinic or PushRelabel solver instance "D" initialized with N nodes
    // and the original graph edges added with their respective capacities.
    void build(const vector<pair<pair<int, int>, long long>>& edges) {
        vector<int> p(n, 0);
        for (int i = 1; i < n; ++i) {
            Dinic D(n, i, p[i]); // Replace with your standard Dinic(N, S, T)
            for (auto& e : edges) {
                D.add_edge(e.first.first, e.first.second, e.second);
                D.add_edge(e.first.second, e.first.first, e.second); // Undirected
            }
            long long flow = D.max_flow();
            tree.push_back({i, p[i], flow});
            for (int j = i + 1; j < n; ++j) {
                if (p[j] == p[i] && D.level[j] != -1) { // j is on i's side of the min-cut
                    p[j] = i;
                }
            }
        }
    }
};
\end{lstlisting}

\begin{itemize}
    \item Initialize via \lstinline|GomoryHu gh(n)| ($0$-indexed). Pass all graph edges to \lstinline|gh.build(edges)|.
    \item Assumes your standard \lstinline|Dinic| structure can report the min-cut partition via reachable nodes (\lstinline|D.level[j] != -1| after \lstinline|max_flow|).
    \item The output \lstinline|tree| contains exactly $N-1$ edges. The all-pairs max-flow/min-cut between any $u$ and $v$ equals the minimum edge weight along the path connecting them in this resulting tree.
\end{itemize}
